%-------------------------
% Resume in Latex
% Author : Aras Gungore
% License : MIT
%------------------------

\documentclass[letterpaper,11pt]{article}

\usepackage{latexsym}
\usepackage[empty]{fullpage}
\usepackage{titlesec}
\usepackage{marvosym}
\usepackage[usenames,dvipsnames]{color}
\usepackage{verbatim}
\usepackage{enumitem}
\usepackage{cmap}
\usepackage[T2A]{fontenc}
\usepackage[utf8]{inputenc}
\usepackage[hidelinks]{hyperref}
\usepackage{fancyhdr}
\usepackage[english, russian]{babel}
\usepackage{tabularx}
\usepackage{hyphenat}
\usepackage{fontawesome}
\input{glyphtounicode}


%---------- FONT OPTIONS ----------
% sans-serif
% \usepackage[sfdefault]{FiraSans}
% \usepackage[sfdefault]{roboto}
% \usepackage[sfdefault]{noto-sans}
% \usepackage[default]{sourcesanspro}

% serif
% \usepackage{CormorantGaramond}
% \usepackage{charter}


\pagestyle{fancy}
\fancyhf{} % clear all header and footer fields
\fancyfoot{}
\renewcommand{\headrulewidth}{0pt}
\renewcommand{\footrulewidth}{0pt}

% Adjust margins
\addtolength{\oddsidemargin}{-0.5in}
\addtolength{\evensidemargin}{-0.5in}
\addtolength{\textwidth}{1in}
\addtolength{\topmargin}{-.5in}
\addtolength{\textheight}{1.0in}

\urlstyle{same}

\raggedbottom
\raggedright
\setlength{\tabcolsep}{0in}

% Sections formatting
\titleformat{\section}{
  \vspace{-6pt}\scshape\raggedright\large
}{}{0em}{\underline}[]

% Ensure that generate pdf is machine readable/ATS parsable
\ifdefined\pdfgentounicode
  \pdfgentounicode=1
\fi

%-------------------------
% Custom commands

\newcommand{\resumeItem}[1]{
  \item\small{
    {#1 \vspace{-2pt}}
  }
}


\newcommand{\resumeSubheading}[4]{
  \vspace{-2pt}\item
    \begin{tabular*}{0.97\textwidth}[t]{l@{\extracolsep{\fill}}r}
      \textbf{#1} & #2 \\
      \textit{\small#3} & \textit{\small #4} \\
    \end{tabular*}\vspace{-7pt}
}


\newcommand{\resumeSubSubheading}[2]{
    \vspace{-2pt}\item
    \begin{tabular*}{0.97\textwidth}{l@{\extracolsep{\fill}}r}
      \textit{\small#1} & \textit{\small #2} \\
    \end{tabular*}\vspace{-7pt}
}


\newcommand{\resumeEducationHeading}[6]{
  \vspace{-2pt}\item
    \begin{tabular*}{0.97\textwidth}[t]{l@{\extracolsep{\fill}}r}
      #1 & #2 \\
      #3 & #4 \\
      #5 & #6 \\
    \end{tabular*}\vspace{-5pt}
}


\newcommand{\resumeProjectHeading}[2]{
    \vspace{-2pt}\item
    \begin{tabular*}{0.97\textwidth}{l@{\extracolsep{\fill}}r}
      \small#1 & #2 \\
    \end{tabular*}\vspace{-7pt}
}


\newcommand{\resumeOrganizationHeading}[4]{
  \vspace{-2pt}\item
    \begin{tabular*}{0.97\textwidth}[t]{l@{\extracolsep{\fill}}r}
      \textbf{#1} & \textit{\small #2} \\
      \textit{\small#3}
    \end{tabular*}\vspace{-7pt}
}

\newcommand{\resumeSubItem}[1]{\resumeItem{#1}\vspace{-4pt}}

\renewcommand\labelitemii{$\vcenter{\hbox{\tiny$\bullet$}}$}

\newcommand{\resumeSubHeadingListStart}{\begin{itemize}[leftmargin=0.15in, label={}, topsep=0pt, itemsep=0pt, parsep=0pt, partopsep=0pt]}
\newcommand{\resumeSubHeadingListEnd}{\end{itemize}}
\newcommand{\resumeBulletListStart}{\begin{itemize}[leftmargin=0.22in, label={--}, topsep=2pt, itemsep=2pt, parsep=0pt, partopsep=0pt]}
\newcommand{\resumeBulletListEnd}{\end{itemize}}
\newcommand{\resumeItemListStart}{\begin{itemize}}
\newcommand{\resumeItemListEnd}{\end{itemize}\vspace{-5pt}}

%-------------------------------------------
%%%%%%  RESUME STARTS HERE  %%%%%%%%%%%%%%%%%%%%%%%%%%%%


\begin{document}

%---------- HEADING ----------

\begin{center}
    \textbf{\scshape Георгий Концевик | Геоаналитик} \\ \vspace{3pt}
    \small
    \hspace{.5pt} {+7 911 846 51 76}
    $|$
    \hspace{.5pt} \href{mailto:george.kontsevik@gmail.com}{george.kontsevik@gmail.com}
    $|$
    \hspace{.5pt} \href{https://github.com/GeorgeKontsevik}{GitHub}
    $|$
    \hspace{.5pt} \href{https://scholar.google.com/citations?user=S--bTg8AAAAJ&hl=ru}{Scholar}
    $|$
    \hspace{.5pt} \href{https://www.linkedin.com/in/george-kontsevik/}{LinkedIn}
    \\ \vspace{3pt}
    \hspace{.5pt} {Санкт-Петербург, Россия}
\end{center}



\section{О себе}
  \vspace{2pt}
  \resumeBulletListStart
    \item Геоаналитик уровня Middle+ с 5 годами опыта.
    \item Перевожу сложные данные в прикладные решения для бизнеса и города.
    \item Специализируюсь на пространственном анализе городской и транспортной инфраструктуры и графовых моделях.
    \item Повышал доступность сервисов, максимизировал входящие потоки, оптимизировал маршруты.
    \item Работаю на Python с библиотеками OSMnx, PySAL, Networkx; владею SQL, Git, Linux.
  \resumeBulletListEnd



%----------- SKILLS -----------

\section{Навыки}
  \vspace{2pt}
  \resumeSubHeadingListStart
    \item Python, SQL, Git, GeoPandas, Shapely, OSMnx, Networkx, PySAL, PostGIS, Linux.
  \resumeSubHeadingListEnd



%----------- EXPERIENCE -----------

\section{Опыт работы}
  \vspace{2pt}
  \resumeSubHeadingListStart
  \resumeSubheading
      {Геоаналитик | Институт дизайна и урбанистики ИТМО}{Сентябрь 2021 \textbf{--} Июль 2026}
      {}{}

  R\&D-лаборатория по разработке IT-решений для городской среды, исполнительный партнер Сбера, OTS Lab, MLA+, Urbanica.

  \resumeSubheading
      {}{}{}{}
  \resumeSubHeadingListEnd
  \resumeBulletListStart
      \item Определил приоритеты развития транспортной сети и размещения логистических хабов для повышения доступности сервисов и эффективности логистики на основе анализа 32 пространственных графов с учетом транспортной структуры и климатических факторов.
      \item Ускорил построение маршрутов до 80 раз при управляемой точности за счёт кластеризации дорожного графа и сокращения области поиска кратчайшего пути.
      \item Повысил точность прогнозирования миграционных потоков, восстановив неполный граф связей между населенными пунктами с помощью гравитационной модели.
      \item Автоматизировал подбор характеристик узлов сети для максимизации входящих потоков, обучив CatBoost-модель на параметрах качества жизни населенных пунктов.
      \item Оценил численность населения по каждому дому города, восстановив пропущенные данные о высотности зданий с точностью до 1 этажа и ошибкой менее 15\% с помощью дерева решений и пространственной автокорреляции.
      \item Ускорил обработку жалоб горожан в соцсетях за счёт разработки алгоритма извлечения адресов из неструктурированного текста и перехода с платного геокодера на бесплатный.
  \resumeBulletListEnd

%----------- EDUCATION -----------

\section{Научно-преподавательская деятельность}
  \vspace{2pt}
  \resumeBulletListStart
    \item Разработал и провел 3 потока магистерского курса по пространственному анализу для студентов-урбанистов на английском языке. Научил с нуля обрабатывать сырые данные и проводить пространственный анализ и моделирование.
    \item Участник международных научных конференций, включая конференции высшего уровня A* в IT сфере (AAAI'26 AI for Urban Planning Workshop).
    \item Автор 12 научных публикаций, из них 3 — в журналах Q1 и 2 — в журналах Q2. \href{https://scholar.google.com/citations?user=S--bTg8AAAAJ&hl=ru}{Ссылка на Scholar.}
  \resumeBulletListEnd

\section{Образование}
  \vspace{3pt}
  \resumeSubHeadingListStart
    
  \resumeEducationHeading
   {Университет ИТМО}
   {}
   {Аспирантура: Управление в организационных системах}{2026}
   {Магистратура: Прикладная информатика}  {2023}
    
  \resumeSubHeadingListEnd

        
\end{document}
